\section{Alcance de la Implementación}

\subsection{Alcance temático}

El proyecto comprende la implementación y parametrización de los siguientes módulos de \textbf{Odoo 19} en modalidad nube, utilizando exclusivamente funcionalidades estándar:

\begin{itemize}
    \item \textbf{CRM Turístico:} Centralización de la base de datos de contactos y pasajeros (PAX), con almacenamiento seguro de documentos adjuntos (fotos de pasaportes) y registro de notas internas e interacciones.
    \item \textbf{Módulo de Proyectos (Tablero Kanban):} Control visual de los servicios y reservas mediante un tablero con etapas operativas: Borrador de Cotización, Confirmación (Adelanto), Compra de Tickets/Hoteles y Ejecución.
    \item \textbf{Gestión de Actividades y Alertas:} Programación de recordatorios manuales (llamadas, correos, tareas) y visualización en la vista de calendario nativa.
    \item \textbf{Gestión de Cobros y Saldos (Facturación básica):} Registro manual de adelantos y saldos recibidos, generación de comprobantes o recibos internos y visualización del estado de pago asociado a cada reserva.
\end{itemize}

Quedan fuera de esta fase (alcance excluido):

\begin{itemize}
    \item Integración automática vía API con portales web de proveedores turísticos.
    \item Implementación de contabilidad financiera completa o facturación electrónica integrada.
    \item Desarrollo de portales de autogestión para el cliente final (Extranet).
    \item Migración histórica de chats o correos electrónicos antiguos.
    \item Desarrollo de reportes de inteligencia de negocios (BI) fuera de las vistas estándar.
    \item Implementación de sitio web, e--commerce o pasarelas de pago.
    \item Gestión de inventarios físicos o activos fijos.
\end{itemize}

\subsection{Alcance organizacional}

La implementación impacta principalmente en:

\begin{itemize}
    \item \textbf{Dirección / Gerencia:} Visión consolidada de clientes, reservas y estado de pagos.
    \item \textbf{Equipo Comercial:} Gestión de oportunidades, comunicación con prospectos y registro estructurado de la información del PAX.
    \item \textbf{Operaciones / Reservas:} Seguimiento de la prestación del servicio mediante el tablero de proyectos.
    \item \textbf{Administración:} Verificación de adelantos y saldos antes de comprometer compras no reembolsables.
\end{itemize}

\subsection{Alcance temporal}

El proyecto se ejecutó entre el 15 de diciembre de 2025 y el 5 de marzo de 2026, estructurado en tres sprints principales:

\begin{itemize}
    \item \textbf{Sprint 1:} Core \& CRM (19 Ene -- 01 Feb).
    \item \textbf{Sprint 2:} Operaciones (02 Feb -- 15 Feb).
    \item \textbf{Sprint 3:} Finanzas \& QA (16 Feb -- 28 Feb).
\end{itemize}

\subsection{Cronograma y Estado de Implementación}

\begin{table}[h]
\centering
\renewcommand{\arraystretch}{1.5}
\begin{tabular}{|l|p{7cm}|p{5cm}|}
\hline
\textbf{Sprint} & \textbf{Objetivo del Sprint} & \textbf{Estado / Entregables} \\ \hline
1: Core \& CRM & Establecer la base de datos y visión 360° del PAX. & Instancia activa, usuarios y permisos configurados. CRM operativo con repositorio de documentos para pasaportes. \\ \hline
2: Operaciones & Digitalizar el flujo logístico y de reservas. & Módulo de Proyectos operativo con tablero Kanban y etapas Borrador, Confirmación, Compra y Ejecución. Actividades y alertas configuradas. \\ \hline
3: Finanzas \& QA & Control de saldos, capacitación y cierre. & Facturación básica configurada para registro de adelantos. Usuarios clave capacitados. Flujo de punta a punta probado y validado. \\ \hline
\end{tabular}
\caption{Cronograma y estado de implementación del proyecto Machu Picchu Exclusive Tours.}
\end{table}

